\documentclass[fleqn]{article}
\usepackage[utf8]{inputenc}
\usepackage[T2A]{fontenc}
\usepackage[russian]{babel}
\usepackage{graphicx}
\usepackage{amsmath}
\renewcommand{\arraystretch}{1.5}
\renewcommand{\baselinestretch}{1.75}
\usepackage[margin=3cm]{geometry}
\usepackage{graphicx} % основной пакет для работы с графикой

\usepackage{amssymb}
\usepackage{caption}
\setlength{\mathindent}{0pt}

\begin{document}

\title{Задачка}
\author{Цуканников Кирилл Андреевич}
\maketitle

\newpage
Числа p и qслучайно выбраны на отрезках: $p \in [2, 6], q \in [0, 4] $ Найти вероятность того, что корни уравнения $x^2 + px + q = 0 $ действительными и различными. \\
Решение: \\
$ D = p^2 - 4q > 0 $ (2 корня строго) \\
$ q < \frac{p^2}{4} $ \\

\begin{figure}[ht]
    \centering
    \includegraphics[width=0.5\textwidth]{2.png}
    \caption{ось обцисс - p, ординат - q}
    \label{fig:my_image}
\end{figure}

Площадь прямоугольника  $p \in [2, 6], q \in [0, 4] $ будет равно 16 \\
Найду площадь $ q < \frac{p^2}{4} $в этом прямоугольнике. Эту площадь найду как сумму прямоугольника с площадью 2 * 4 = 8 и фигуры под $q = \frac{p^2}{4}$ \\
$\int\limits^4_2 \frac{p^2}{4} dp  = \frac{14}{3}$
$S = \frac{14}{3} + 8 = \frac{38}{3}$(площадь под условием $ q < \frac{p^2}{4} $) \\
Теперь поделю площадь удовлетворяющую условию га плозадь общую прямоугольника. То есть вопсопльзуюсь геомтерическим определением вероятности \\
$ \frac{\frac{38}{3}}{16} = \frac{19}{24} $ - ответ.
\end{document}